% ==============================================================================
% Fichier : chapitres/07_evaluation_tests.tex
% Livre VII : Évaluation et Tests de la Sécurité
% ==============================================================================

\chapter{Évaluation et Tests de la Sécurité}
\label{chap:tests_securite}

\begin{boxobjectifs}
\textbf{Objectifs du Livre~VII :} Comprendre l'importance des tests de sécurité, maîtriser leurs types (boîte noire/grise/blanche) et les différentes étapes d'un test d'intrusion (autorisation, reconnaissance, scanning, exploitation, élévation, nettoyage, reporting).
\end{boxobjectifs}

% ------------------------------------------------------------------------------
\section{L'Importance des Tests de Sécurité}
% ------------------------------------------------------------------------------

Les tests de sécurité permettent d'identifier proactivement les vulnérabilités avant qu'un attaquant malveillant ne les exploite. Ils constituent un pilier fondamental de la démarche de sécurité défensive et sont souvent requis par les référentiels (ISO~27001, PCI~DSS).

% ------------------------------------------------------------------------------
\section{Types de Tests de Sécurité}
% ------------------------------------------------------------------------------

\begin{table}[H]
\centering
\small
\begin{tabular}{p{3cm}p{5.5cm}p{4cm}}
    \toprule
    \textbf{Type} & \textbf{Description} & \textbf{Usage} \\
    \midrule
    \textbf{Boîte Noire} & L'auditeur n'a aucune information sur la cible. Simule un attaquant externe. & Évaluation de l'exposition externe. \\
    \textbf{Boîte Grise}  & L'auditeur dispose d'informations partielles (ex. : compte standard). & Simulation d'un utilisateur interne. \\
    \textbf{Boîte Blanche} & Accès complet au code source et à l'architecture. & Audit de code, revue de sécurité approfondie. \\
    \textbf{Red Team}     & Équipe attaquante simulant une APT (menace persistante avancée). & Test de la capacité de détection et de réponse. \\
    \textbf{Blue Team}    & Équipe défensive (SOC) répondant aux attaques de la Red Team. & Amélioration des défenses et de la réactivité. \\
    \bottomrule
\end{tabular}
\caption{Types de tests de sécurité}
\end{table}

% ------------------------------------------------------------------------------
\section{Les Étapes d'un Test de Sécurité (Pentest)}
% ------------------------------------------------------------------------------

\subsection{a) Autorisation et Périmètre}
Étape préalable indispensable. Sans autorisation écrite (ROE — Rules of Engagement), toute intrusion est illégale. Définir : les systèmes inclus, les systèmes exclus, les actions autorisées, les horaires, les contacts en cas d'incident.

\subsection{b) Reconnaissance (Passive)}
Collecte d'informations sans contact direct avec la cible. Outils OSINT :
\begin{itemize}
    \item \textbf{Maltego :} Cartographie des relations (domaines, emails, organisations).
    \item \textbf{Shodan :} Moteur de recherche d'équipements connectés à Internet.
    \item \textbf{theHarvester :} Collecte d'emails, sous-domaines, IPs via moteurs de recherche.
    \item \textbf{WHOIS, dig, nslookup :} Informations DNS et d'enregistrement de domaine.
\end{itemize}

\subsection{c) Scanning et Énumération (Active)}
Interaction directe avec la cible pour cartographier les ports, services et systèmes.
\begin{itemize}
    \item \textbf{Nmap :} Scanner de ports et de systèmes d'exploitation.
    \item \textbf{Nessus / OpenVAS :} Scanners de vulnérabilités automatisés.
    \item \textbf{Angry IP Scanner, Advanced IP Scanner :} Découverte d'hôtes actifs sur un LAN.
    \item \textbf{Dirb / Gobuster :} Découverte de répertoires cachés sur les serveurs web.
\end{itemize}

\begin{lstlisting}[language=bash, caption=Commandes Nmap courantes]
# Scan de ports SYN (furtif)
nmap -sS -p 1-1024 192.168.1.0/24
# Détection de version et OS
nmap -sV -O 192.168.1.50
# Scan complet avec scripts par défaut
nmap -A 192.168.1.50
\end{lstlisting}

\subsection{d) Recherche des Vulnérabilités}
Analyse des résultats de scanning pour identifier les failles exploitables.
\begin{itemize}
    \item \textbf{CVE (Common Vulnerabilities and Exposures) :} Base de données des vulnérabilités publiques.
    \item \textbf{ExploitDB / SearchSploit :} Base d'exploits publics.
    \item \textbf{Rapid7 / Metasploit :} Référentiel d'exploits fonctionnels.
    \item \textbf{Nikto :} Scanner de vulnérabilités des serveurs web.
    \item \textbf{CVSS (Common Vulnerability Scoring System) :} Score de gravité de 0 à 10.
\end{itemize}

\subsection{e) L'Exploitation}
Tentative de compromission effective via les vulnérabilités identifiées.
\begin{itemize}
    \item \textbf{Metasploit Framework :} Standard de l'exploitation de vulnérabilités système.
    \item \textbf{BurpSuite / OWASP ZAP :} Proxy d'interception et d'attaque des applications web.
    \item \textbf{SQLMap :} Automatisation des injections SQL.
    \item \textbf{Mimikatz :} Extraction de credentials depuis la mémoire Windows (LSASS).
\end{itemize}

\subsection{f) L'Élévation de Privilèges}
Après l'accès initial, l'auditeur cherche à obtenir des droits administrateur/root pour prouver l'impact maximal.
\begin{itemize}
    \item \textbf{Linux :} Exploitation SUID, sudo misconfiguration, kernel exploits.
    \item \textbf{Windows :} Token impersonation, unquoted service path, AlwaysInstallElevated.
\end{itemize}

\subsection{g) Le Nettoyage}
Suppression de toutes les traces laissées (fichiers déposés, comptes créés, règles pare-feu modifiées) pour remettre le système dans son état initial.

\subsection{h) Le Reporting}
Rédaction du rapport de pentest incluant :
\begin{enumerate}
    \item \textbf{Synthèse exécutive} (pour la direction).
    \item \textbf{Méthodologie} et périmètre testé.
    \item \textbf{Findings} détaillés (titre, score CVSS, description, PoC, impact, remédiation).
    \item \textbf{Plan d'action} priorisé.
\end{enumerate}

% ------------------------------------------------------------------------------
\section{Évaluation des Acquis --- Livre~VII}
% ------------------------------------------------------------------------------

\begin{boxexo}
\textbf{QCM :}
\begin{enumerate}
    \item Un test boîte noire simule : \textbf{(a) Un auditeur interne avec accès complet \quad (b) Un attaquant externe sans information préalable \quad (c) Un utilisateur avec un compte standard \quad (d) Un développeur avec le code source}
    \item Nmap est un outil de : \textbf{(a) Chiffrement \quad (b) Scanning de ports \quad (c) Gestion des identités \quad (d) Journalisation}
    \item Le CVE est : \textbf{(a) Un outil d'exploitation \quad (b) Un dictionnaire public des vulnérabilités \quad (c) Un protocole réseau \quad (d) Un modèle de contrôle d'accès}
    \item Le nettoyage en pentest consiste à : \textbf{(a) Formater les disques \quad (b) Supprimer les traces laissées pendant le test \quad (c) Réinstaller le système \quad (d) Désactiver les pare-feux}
\end{enumerate}

\textbf{Travail Pratique (ENSPY) :}\\
Réaliser une évaluation de la sécurité de l'application de gestion des notes de l'ENSPY. Identifier les vulnérabilités (OWASP Top 10), les scorer avec CVSS et proposer des remédiations.
\end{boxexo}
